\section{Data and Preprocessing}\label{sec:data_preprocessing}
The section will detail what data was used for training the different models, 
as well as how it was preprocessed, with the purpose of enabling 
reproducibility. All the code used for the project is available on 
GitHub\footnote{link to github}. For transparency it should be noted that 
generative AI tools have been used extensively for both writing and reviewing 
the code. 
\subsection{Data Sources}\label{sec:data_sources}
\subsubsection{Weather data}
The weather data used in the paper originates from \gls{dmi}. They collect 
weather data from their numerous weather stations located in Denmark, 
Greenland, and the Faroe Islands. They then provide weather forecasts on a 
national scale. \\
They publish their collected data through both API and downloadable files. They 
also provide the data in different granularities, such as by station, 
municipality, country, or by 10/20 square kilometre grid 
areas\footnote{https://opendataapi.dmi.dk/v2/climateData/bulk/}. For this 
paper, the data was accessed on a granularity level of individual stations. 
Additionally, \gls{dmi} provides two different versions of meteorological 
measurement data, one 
version that is purely raw data, and one that has been processed for quality 
control by \gls{dmi} employees. The quality controlled version of the data was 
chosen for this paper. The utilized data has a date range from 00:00:00 - 
01/01/2015 until 00:00:00 - 
03/25/2026, the end of the range signifies the date the download was performed. 
Downloading the data was preferred to avoid multiple API calls, as there were a 
limit of $300.000$ measurements accessible through the API in each call. If one 
wishes to develop a prediction model utilizing this data, and update it during 
deployment, API calls would be the preferred method, to get limited daily data 
amounts. 

\subsubsection{Energy production \& consumption data}
The data regarding energy \gls{pc} originates from 
\gls{en}. They are a subsidiary of the Danish government and manages the Danish 
energy grid. The also publish a variety of data regarding the Danish energy 
grid, again, through both API and downloadable files. For the same reason as 
with \gls{dmi}, downloadable files were chosen. The \gls{pc} data used for this 
paper is acquired from the "Production and Consumption - Settlement" 
publication\footnote{\url{https://www.energidataservice.dk/tso-electricity/productionconsumptionsettlement}}.

\subsubsection{Price data}
The price data also originates from \gls{en}. In October of 2025 \gls{en} 
changed their reporting format to follow the new 15-minute timing practices of 
Nord Pool. Due to this, the previous publication was discontinued, and a new 
one 
was created in accordance with the new format. The price data for this paper 
has been concatenated between these two different publications: Elspot 
Prices\footnote{https://www.energidataservice.dk/tso-electricity/Elspotprices}, 
from before October 2025, and 
Day-Ahead-Prices\footnote{https://www.energidataservice.dk/tso-electricity/DayAheadPrices},
for after October 2025. To ensure compatibility between these two formats, we 
simply discard prices that are not exactly hourly in the newer dataset.

\subsubsection{Midas Energy data}
Midas Energy has also provided weather data in the format they use for their 
own prediction models. This is, in simple terms, a processed and generalised 
version of the \gls{dmi} data. This data will later be used as a substitution 
for the weather data, to gauge the difference between the two varieties of 
weather data. It will also be added to the best performing models, to 
investigate its direct impact. 

\subsection{Data format and Preprocessing}
This section details the data formats before and after preprocessing, as well 
as the preprocessing pipeline for each data source.
\begin{table*}[t]
	\centering
	\caption{Snapshot of a Data Record (Raw vs. 
		Preprocessed)}
	\label{tab:data_snapshot_wide}
	\begin{tabular}{@{} l l l l l l l l @{}}
		\toprule
		\multicolumn{8}{@{}l}{\textbf{Raw Data 
				Snapshot}} \\
		\midrule
		\textbf{Timestamp} & \textbf{Zone} & 
		\textbf{Spot Price} & 
		\textbf{Load} & \textbf{Wind Prod.} & 
		\textbf{Solar Prod.} & 
		\textbf{Temperature} & 
		\textbf{Precipitation} \\
		2026-04-26 12:00   & DK1           & 45.20 
		EUR           & 3200 
		MW       & 12.5 MW             & 400 
		MW               & 14.2\textdegree 
		C    & 0.0 mm                 \\
		
		\midrule
		\midrule
		
		\multicolumn{8}{@{}l}{\textbf{Preprocessed 
				Data Snapshot}} \\
		\midrule
		\textbf{Timestamp} & \textbf{Zone\_DK1} & 
		\textbf{Price\_Norm} & 
		\textbf{Load\_Norm} & \textbf{Wind\_Norm} 
		& \textbf{Solar\_Norm} & 
		\textbf{Temp\_Norm}  & 
		\textbf{Precip\_Flag} \\
		1761476400         & 1                  & 
		0.23                 & 
		0.65                & 0.41                
		& 0.12                 & 
		0.55                 & 
		0                     \\
		\bottomrule
	\end{tabular}
	\caption{\textbf{PLACEHOLDER - NOT REAL DATA DEPICTED.}}
\end{table*}
\subsubsection{Weather data}
\paragraph{Raw data}
The raw weather data as can be seen partly in table X, consists of various 
real-time measurements performed by a variety of instruments located at 
different weather stations. Each row in the raw data consists of: a measurement 
ID, the geographical coordinates for the measurement, the measurement parameter 
type, the time resolution with the related timeslots, the value, and the 
station ID. It also includes several features related to \gls{dmi} quality 
assurance, such as creation time, calculation time, quality control status, 
etc., however, these are discarded as they have no metrological value. 
Further, there exists four different timestamps within each entry. For the 
purposes of this paper, we utilise the 'from' timestamp as the main timestamp, 
as that will correspond to the price from 14:00 to 15:00.
\paragraph{Preprocessed data}
Preprocessing on the weather data mainly consisted of grouping each individual 
measurement by station ID, then aggregating and averaging measurements to 
become hourly if they are not already. Lastly, we once again aggregate these 
hourly measurements by region, using both station ID as well as measurement 
coordinates, to determine the region, ending up with one hourly dataset for the 
DK1 and DK2 regions respectively. During this geographical segregation we also 
excluded both Greenland, the Faroe Islands, and Bornholm, due to the 
geographical distance from these regions to DK1 and DK2. This resulted in data 
such as the example depicted in table X.
Further validation of the preprocessed weather data was performed, in order to 
ensure data quality. Features that did not have data for the entire period were 
also entirely removed. This means that if a type of sensor was decommissioned 
during the decade of training data its feature data would be entirely excluded 
from training, to enable rolling-window cross-validation, as will be detailed 
in \autoref{sec:experiments}.

\subsubsection{Production and Consumption data}
The \gls{pc} data required minimal preprocessing, as it was already published 
in the hourly format. The only preprocessing performed on this dataset was 
dividing it into the two regions, DK1 and DK2, which was already labelled 
within a column of the dataset. Then we ensured that the date-range matched the 
\gls{dmi} data. The data contained within the dataset consists of hourly 
production of wind energy of various size-categories of wind-farms, both off- 
and onshore, solar energy production of various size-categories of solar farms, 
fossil-fuel related energy production, exchanges with neighbouring countries, 
as well as Danish energy consumption data. Notably, columns relating to foreign 
energy exchange contained frequent missing values when no exchange had occurred 
between the countries during that specific hour. Since these missing values 
correspond to an energy exchange of zero, they were explicitly encoded as such. 
This ensures the models accurately interpret periods of national grid 
isolation, without the need for interpolation.

\subsubsection{Price}
The price data was divided into two publication, due to the previously 
mentioned Nord Pool formatting change. However, the columns in the tables where 
still the same, so we just ensured that they were in the same order, and then 
appended the most recent publication to the earlier one. We also discarded 
non-hourly prices to keep formatting consistent across tables. 

\subsubsection{General Preprocessing}
Additionally, general preprocessing was performed across all datasets, or to 
ensure consistent UTC time formatting, also removing the issue of daylight 
savings time anomalies. Lastly, all datasets were sorted into ascending 
timestamps, and data-ranges were enforced over all datasets.

\subsubsection{Master Matrix Data}
After preprocessing, all of these different tables were joined by their 
timestamp, creating one master table, to be used for model training. 
Additionally, several columns were created for the express purpose of either 
enhancing model performance or combatting over-fitting of the models during 
training. Four different categories of columns were added: 
\paragraph{Date \& Time}
Not to be confused with date-time, or timestamps, these columns instead relate 
to representing the hour of the day, the day of the week, the month, 
and the day of the year as ascending integers. We add these to give the models 
an integer representation of the days and hours, to better let them interpret 
consumption patterns, e.g. the differing patterns present during weekdays and 
weekends. To enable detection of such patterns further, we also encode these 
integers as Sine and Cosine waves, creating a scaled version of them, that may 
be more useful to certain model types. Utilizing both mitigates the natural 
symmetrical ambiguity that would be present otherwise, mapping each time step 
to a unique coordinate on a unit circle. The waves were created using the 
following formulae \ref{eq:sine_enc} and \ref{eq:cos_enc}:
\begin{align}
	x_{\text{sin}} &= \sin\left(\frac{2\pi x}{\max(x)}\right) 
	\label{eq:sine_enc} \\
	x_{\text{cos}} &= \cos\left(\frac{2\pi x}{\max(x)}\right) \label{eq:cos_enc}
\end{align}
where $x$ is the integer being encoded, and $\max(x)$ is the maximum value of 
that specific period e.g., 24 for the hour of the day, or 365 for the day of 
the year. Notably, day of month was deliberately excluded as a feature, 
because of its variable period length, due to not all months contain the same 
number of days, would introduce gaps at the end of shorter months that cannot 
be cleanly resolved with a static denominator. A dynamic denominator, adjusting 
$\max(x)$ to reflect the actual number of days in each specific month, would 
resolve this, however, the same information is implicitly captured by the 
combination of the day of year and month encodings already present.

\paragraph{T-Minus data}
Since we will be training tree-based models, which does not possess any 
inherent memory capabilities, we introduce T-Minus data, or lagging data, which 
are essentially just copies of earlier data at various horizons. All 
weather-related measurements, as well as \gls{pc} related values, have been 
copied and inserted 24-hours later, giving the model inherent access to 
yesterdays weather and consumption at this exact hour, in order to predict the 
price for the current hour. Additionally, for the price data specifically, 
t-minus data was included for the previous 24-hour, 48-hour, and 168-hour, as 
previous prices were deemed especially relevant for predicting the current 
price.
\paragraph{Targets}
Originally, the absolute price was used as the target for model training. 
However, to combat overfitting, a new target was introduced to the master 
matrix: the delta target. Instead of trying to predict the absolute price, the 
models were instructed to predict the change in price.
\begin{equation}\label{eq:delta}
	\Delta P_{t,h} = P_{t+h} - P_t
\end{equation}
where $P_t$ is the price at the current hour and $h$ is the forecast 
horizon in hours.\\
The results of 
both training approaches will be discussed in \autoref{sec:results}.
\paragraph{Groupings}
To make training different variations of the data easier, each column was 
assigned a group consisting of related columns. These groupings will be 
discussed further in \autoref{sec:experiments}.

\paragraph{Missing data and imputation}
To ensure satisfactory data quality, features that are more than 90\% 
empty, or have constant values in 95\% of rows, are dropped entirely prior to 
training. These thresholds were chosen conservatively to retain as many 
features as possible, under the assumption that moderately sparse features 
could be recovered through imputation. A post-hoc audit confirmed that no 
columns were dropped under the emptiness criterion across any dataset. The 
constant value criterion resulted in the removal of two internal DMI 
quality control columns, \texttt{Stations\_Reporting\_Min} and 
\texttt{Stations\_Reporting\_Max}. These columns carry no 
meteorological value and their removal is appropriate.

For remaining missing values, three steps were performed to all columns. First, 
interpolation is applied to gaps of up to 12 consecutive hours, to fill in 
smaller gaps. Second, forward and backward filling are applied as a 
supplementary pass, also capped at 12 hours each, to handle edge cases where 
interpolation is insufficient, due to a one of the sides of the gap being open. 
Any values still missing after both passes, implying a consecutive gap 
exceeding 24 hours, are filled using the column mean as a last resort.

The behaviour of this strategy differed substantially between data 
sources. The price data required no imputation whatsoever, containing 
zero missing values across both publications. The weather data was 
largely clean, with 0.25\% and 0.53\% of values originally missing 
in DK1 and DK2 respectively. The vast majority of these were 
recovered by interpolation, with the mean fallback triggered only for 
\texttt{temp\_soil\_30} and \texttt{leaf\_moisture} in specific years, 
accounting for at most 0.02\% and 0.18\% of all weather values in DK1 
and DK2 respectively. 

The production and consumption data presented a different situation. 
Approximately 9.87\% of its numeric values were 
originally missing, concentrated entirely in four columns: 
\texttt{UnknownProdMWh}, \texttt{CommercialPowerMWh}, 
\texttt{GridLossInterconnectorsMWh}, and 
\texttt{GridLossDistributionMWh}. These represent production and grid 
measurement categories that were introduced or changed during the 
11-year dataset period. In these cases, interpolation is not physically 
meaningful, and the column mean effectively encodes the absence of a 
measurement category as a neutral baseline value. Additionally, six exchange 
columns containing frequent missing values were zero-filled separately, as a 
missing exchange value physically corresponds to zero energy flow between 
countries rather than an unknown measurement.

We also create columns of binary flags to signal what data is natural and what 
is synthetic. These flags are primarily motivated by 
tree-based models, which lack any inherent mechanism to discount 
uncertain inputs. By explicitly encoding which values are synthetic, tree-based 
models may learn to route predictions away from relying on such criteria when 
other features may be more reliable. Neural network-based 
models can achieve a similar effect through their attention and gating 
mechanisms without explicit flags, however, the flags are nonetheless present 
within the master matrix for all models types, as they carry no cost for models 
that do not benefit from them.\\

\subsubsection{Midas Energy data}
No preprocessing has been performed on the Midas Energy Data, except for 
ensuring consistent time-formatting in relation to the other datasets. 

%https://www.dmi.dk/friedata/dokumentation/climate-data link to description of 
%weather data fields